% ===== §13 Relational Aesthetics and the Four-Step Model =====
\section{Relational Aesthetics and the Cultivation of the Field}
\label{sec:relational-aesthetics}

\noindent
The preceding parts established a chain of grounding. Relational production requires a field; the construction of a field is an aesthetic rather than an engineering matter; the aesthetic and the ethical are coupled at their root. The chain terminates in a single reframing: the cultivation of the capacity to build a generative field is what the word \emph{aesthetic education} properly names. What remains is to say what that capacity actually is, finely enough that it can be practised. This section supplies the two instruments the rest of the paper will use. The first is a concept, \term{relational aesthetics}, which locates beauty neither in the object nor in the subject but in the coupling event between subjects. The second is a model, a four-step cycle of perception, reception, sharing, and reproduction, which is not a procedure laid over experience from outside but the minimal form a good cycle takes when it is enacted between two people in the course of an ordinary day.

\subsection{Beauty is produced relationally}
\label{sub:beauty-relational}

The history of aesthetics has largely been a quarrel between two parties. One party locates beauty in the object: there are proportions, harmonies, or formal properties that a thing possesses, and to find a thing beautiful is to register a feature already there. The other party locates beauty in the subject: beauty names a state of the perceiver, a pleasure that claims universality without resting on a concept, so that the judgement of taste says more about the structure of our faculties than about the rose. The first party cannot explain why competent perceivers disagree without lapsing into a doctrine of error. The second cannot explain why the judgement feels answerable to the thing, why it can be corrected, why two people leaning toward the same evening light feel they are tracking something rather than comparing private weather.

Each party should be granted its strongest form before the third position is proposed, since the third position earns its place only by explaining what defeats the other two. The objectivist is strongest where beauty seems plainly to reside in the thing, in the proportion that holds whoever looks, the interval that sounds whoever hears, and the strength of the position is its account of answerability, of why the judgement of beauty can be right or wrong, corrected, taught, why it is not a mere report of private pleasure. Its weakness is disagreement: if beauty is in the object, the persistent, competent, irreducible disagreement of fine perceivers must be error, and a doctrine that convicts half its best practitioners of error wherever they disagree has misdescribed what they are doing. The subjectivist is strongest exactly where the objectivist is weak, in the account of disagreement, which becomes no error but the natural variety of differently constituted perceivers, and in the account of why beauty moves us, why it is bound up with pleasure and feeling rather than standing as a cold property. Its weakness is answerability: if beauty is a state of the perceiver, the felt sense that the judgement is answerable to the thing, correctable, shared, becomes an illusion, two people who lean toward the same light only seeming to track a common object while in truth comparing private states. Each party is strong exactly where the other is weak, which is the signature of a missing third term that both have approached from opposite sides and neither has reached.

The framework of this series forces a third position, and forces it not as a compromise between the first two but as a consequence of the relational ontology that the series has carried throughout. If the subject is itself generated in relation, and value is the holonomy accumulated along a relational path, then the aesthetic, which we have argued is the perception of appropriateness in a field, cannot be a property of an isolated object or an isolated subject, because neither an isolated object nor an isolated subject is where appropriateness lives. Appropriateness is a feature of a coupling. The same gesture, the same placement of a thing, accumulates good holonomy in one relational state and bad holonomy in another, and there is no reading of the gesture or the thing alone that fixes which. Beauty, on this account, is the felt registration that a coupling is accumulating in the generative direction. It is an event between, and it is produced there, in the between, when it is produced at all.

The third position explains what defeated the other two, which is the test it must pass to deserve the name of a resolution. It explains answerability, the objectivist's strength and the subjectivist's failure, because the coupling is a real event with a real direction, and a judgement about it can be right or wrong about whether the field accumulates, answerable to the coupling though not to the object alone. And it explains disagreement, the subjectivist's strength and the objectivist's failure, because the coupling is between these two in this state, and a perceiver outside the coupling, or the same coupling in another state, may truly find otherwise without error, the variety being the variety of couplings and states rather than of private weather. The two leaning toward the evening light are neither registering a property of the light nor comparing private pleasures; they are, if the moment is real, a coupling accumulating upon the light, and their sense of tracking something together is exactly right, the something being the coupling event their joint attention produces. Relational aesthetics is thus not a third opinion beside two others but the position that locates beauty where both the others were reaching, in the between that the object-party mistook for the object and the subject-party mistook for the subject.

\claim{\emph{Relational aesthetics.} Beauty is not a property of the object and not a state of the subject. It is a coupling event between subjects, produced relationally. The aesthetic perception of appropriateness is the registration that a relational field is accumulating holonomy in the generative direction; such registration has its proper site in the coupling, not in either pole taken alone.}

This is a strong claim and it should be stated at its strongest, because its strength is what does the work later. It does not say that beauty has a relational dimension, or that beauty is enhanced by being shared. Both of those are weak claims compatible with locating beauty in the object or the subject and then adding a social supplement. The claim here is constitutive. The supplement is the thing. A beauty that has not entered the relational register is not yet a beauty that this account recognises, in the way that a vow not yet spoken to anyone is not yet the kind of vow whose causality the series has studied. This has an immediate consequence that the four-step model will make concrete: a beauty perceived and kept inside one subject is incomplete, and the labour of completing it, of carrying it across into the between, is not optional adornment but the act by which beauty becomes what relational aesthetics says it is.

One clarification protects the claim from a fast objection. To say beauty is produced in the coupling is not to say it is produced by agreement, or by two parties happening to report the same pleasure. Coupling is not consensus. Two people can agree, in words, that a thing is beautiful while no coupling event occurs between them; the agreement is then a pair of private judgements laid side by side, the very side-by-side arrangement the second party to the old quarrel mistakes for relation. What relational aesthetics requires is not concord of report but the actual event in which a field, attended by both, accumulates in the generative direction and is felt to do so. The later section on sharing (\xref{sec:sharing}) is wholly occupied with the difference between these two things, because the difference is where the concept earns its keep.

A second objection is more serious and must be answered, because it concerns the very cases that the tradition of aesthetics took as paradigmatic. Surely, the objection runs, the masterwork is beautiful whether or not anyone is in a coupling with anyone; the painting in the empty gallery, the symphony with no one in the hall, are beautiful in themselves, and a theory that locates beauty in the coupling between subjects cannot account for the beauty of the artwork that no coupling attends. The objection is forceful, and the answer is not to deny the beauty of the artwork but to locate it correctly. The artwork is beautiful, on this account, because it is the occasion and the trace of couplings, made to produce them and bearing the marks of the ones that made it. The painting is not beautiful in the way a proportion is a property of a triangle; it is beautiful in that it reliably occasions a coupling between itself and a perceiver, and between perceivers who attend to it together, a coupling the painter built it to occasion and built out of the couplings of his own making. The beauty we ascribe to the artwork in the empty gallery is the beauty of its disposition to occasion such couplings, its standing readiness to produce the event when a perceiver comes, in the way that we call a thing soluble when no one is dissolving it, ascribing to it the disposition though the event is not now occurring. The artwork is a crystallised invitation to coupling, and its beauty is real, the real disposition to occasion the event in which beauty proper is produced.

This locates the museum correctly within the account and reconciles the paper's relational aesthetics with the tradition it displaces. The museum is not the home of a beauty that contradicts the relational account; it is a place full of crystallised invitations to coupling, objects made to occasion the event between a perceiver and the work and between perceivers who share it. The tradition that located beauty in the museum object was not simply wrong but partial: it perceived the object's real disposition to occasion coupling and mistook the disposition for a property complete in the object, the way one might mistake solubility for a property the salt has in itself rather than a disposition it has toward water. Relational aesthetics corrects the partiality without denying what the tradition perceived, and in doing so it earns the right to move aesthetic education from the museum to the kitchen, since the kitchen is where the coupling the museum object only occasions is actually built, the living event of which the artwork is the crystallised invitation. The artwork invites the coupling; the everyday is where the coupling is lived, and the aesthetic education of the everyday is the cultivation of the capacity to build the events that the museum object could only, in its beautiful stillness, invite.

\subsection{The grounds for a model, and for these four steps}
\label{sub:why-four}

If beauty is produced relationally, then the capacity to participate in its production can be cultivated, and a cultivation needs a structure fine enough to be practised against. The structure proposed here has four movements. They are not chosen for symmetry or for the convenience of exposition. They are the decomposition of a single good cycle into the smallest sequence of distinguishable acts by which such a cycle is enacted between two people, and the claim that there are exactly these four, in this order, is a claim that the paper must earn rather than assume.

The first movement is \term{perception}: the registration of what is occurring in the field, whether it is accumulating or spinning in place, what here would be appropriate. This is the input face, the inferential moment. The second is \term{reception}: the perceived must be let in, allowed to alter the one who perceives, rather than held off at the level of cognition or intercepted by defence. The third is \term{sharing}: the received must be carried back into relation, attended by another, for only there does it become the coupling event that relational aesthetics says beauty is. The fourth is \term{reproduction}: the shared does not dissipate but settles as a raised starting point, so that the next perception begins higher, finer, more able to register what before went unseen.

\claim{The four movements form a cycle that closes without returning to its origin. Perception, reception, sharing, and reproduction return the practitioner to perception, but because reproduction has raised the starting point, what is returned to is the root and not the origin. The cycle is therefore not a wheel but a spiral, and its non-return is exactly the anholonomy that the series has named as the geometric signature of a good cycle.}

This is why the model is not arbitrary and not merely didactic. Its form is the form of the thing it cultivates. A reader who has followed the series will recognise the closing-without-return as the figure that ran through the account of poetry, the figure of coming back to the root yet not to the origin, the Berry phase read as the formal correlate of a cycle that gains rather than conserves. The four-step model instantiates that figure at the most intimate scale, the scale of a single returned gesture between two people. It follows that the model can fail in exactly the ways a good cycle can fail. If reproduction merely repeats the same, the spiral collapses to a wheel and the holonomy goes to zero; the practice has become the bad repetition whose metaphysics \xref{sec:maintenance} examines. If sharing does not occur, the cycle never leaves the single subject and beauty is never produced in the sense that matters; the practice has become the cultivation of private taste, which is the old refuge of an aesthetic education that mistook a secondary contradiction for the principal one. The model is thus not a recipe whose steps one performs in order to succeed. It is a diagnosis of where, in a given relation at a given moment, the accumulation has stalled.

\subsection{The four steps under the generative dialectic}
\label{sub:four-dialectic}

The model does not stand apart from the method announced earlier. The four steps are subject to the same analysis of principal and secondary contradiction, and this is what keeps the model from hardening into a fixed list of things one ought to do. At any moment in a relation, the four steps are not equally weighted; one of them is the principal site, the place where the accumulation is currently stalling and where generative attention should therefore be directed. Which step is principal is not fixed by the model but read from the state of the field, and it migrates as the state changes.

\begin{itemize}
  \item When a relation has grown estranged, the principal site is \emph{sharing}. Beauty is still perceived and even received within each subject, but it no longer crosses into the between; the field has lost the coupling event, and the labour that matters is the carrying-across, however small.
  \item When a relation has grown dull, the principal site is \emph{perception}. The partners still act, still share by habit, but no longer see; the contingency that would refresh the field goes unregistered, and the labour that matters is the recovery of sight.
  \item When a relation has lately been wounded, the principal site is \emph{reception}. Defence is highest after injury, and what is perceived and even offered cannot get in; the labour that matters is the lowering of the gate, which no amount of perceiving or sharing can substitute for.
  \item When a relation has stalled in plenty, the principal site is \emph{reproduction}. Beauty is perceived, received, and shared, but consumed rather than thickened, each sharing a re-enactment of the last; the labour that matters is the turning of consumption back into accumulation.
\end{itemize}

The point of this analysis is not to issue four rules but to refuse the single rule. There is no posture that is always the right one, no step that is always the one to work on, because the right labour is a function of which contradiction is currently principal, and that is something to be discerned in the particular field rather than legislated for fields in general. The discernment is itself an exercise of the trained sensibility the paper is about. To read which step is principal is already to perceive the field, which means the dialectical analysis is not applied to the model from outside but is the first movement of the model turned upon the relation as a whole. The method and the object meet here, and their meeting is one more instance of the form enacting its claim.

The relation between the four-step model and the three beauties set out earlier should be made explicit, since the two articulations of the one capacity might otherwise seem to compete. The three beauties, perception and practice and maintenance, are the anatomy of the capacity, its cross-section, the three faces a field-builder has; the four steps, perception and reception and sharing and reproduction, are the same capacity in motion, the cycle the anatomy turns. The two are not rival lists but the static and dynamic descriptions of one thing, and they map onto each other in a definite way. The beauty of perception is the first movement of the cycle, the perceiving that opens it. The beauty of practice spans the middle movements, reception and sharing, since to receive the contingency and to carry it into the between are the doing in which the apt act is performed, the practice that builds the field. The beauty of maintenance is the fourth movement, reproduction, the sustaining across time that keeps the cycle turning and the field deepening. The anatomy has three faces and the cycle has four movements because the middle of the cycle, the doing, divides into the receiving and the sharing that the static anatomy gathers under the single face of practice, the cycle resolving into four what the anatomy holds as three. Neither articulation is more correct; the anatomy shows the capacity's parts and the cycle shows their turning, and the paper needs both, the cross-section to exhibit what the field-builder has and the cycle to exhibit how what she has is exercised, the two descriptions of the one capacity that the whole paper is devoted to cultivating.

This double articulation is itself a small instance of the paper's method, the holding of two true descriptions without forcing one to subsume the other. The anatomy and the cycle are both true of the capacity, and neither is the description of which the other is an application; to insist that the capacity is really three faces, with the four steps a derivative regrouping, or really four steps, with the three faces a crude simplification, would be to settle into a single governed description the plurality that the paper has everywhere refused to settle. The capacity is both, a thing with three faces and a cycle with four movements, and the two articulations illuminate it from different angles, the static and the dynamic, neither exhausting it. That the one capacity admits two true and non-competing articulations is not a defect of the account but a feature of the thing, which is rich enough to be truly described in more than one way, and the paper's holding of both, without subordinating either, is the same fidelity to plurality that its polyphonic close will perform at the scale of the whole.

\subsection{The work of the following sections}
\label{sub:four-roadmap}

Each of the four movements carries a body of theory, and in two of them the theory is dense enough, and the traditions that converge on them divergent enough, that the movement requires a section of its own. Perception (\xref{sec:perception-field}) gathers predictive processing, the Buddhist account of undivided awareness, and the phenomenology of the everyday, and asks how the field becomes visible at all. Reception (\xref{sec:reception}) gathers the ethics of welcoming the other, the psychoanalytic account of containment, and the Daoist figure of the mirror that meets without storing, and asks what it is to let beauty and the other in. Sharing (\xref{sec:sharing}) gathers joint attention, witnessing, and the dialogical generation of meaning, and asks how the relational is born. Reproduction (\xref{sec:reproduction}) gathers the philosophy of repetition and the political economy of the reproductive act, and asks how a cycle thickens in time without either hoarding or dulling. A fifth section (\xref{sec:neuro}) gives the implementation-level interface, the observable signatures each movement should leave if the account is true, with the restraint that an interface is not an identity. None of these sections is a separate doctrine. Each is one movement of a single cycle, and the cycle is the form of the good cycle itself.
